\begin{afterword}
\summaryhead

The only perfect map of California is California. A useful map leaves things out, and a brain is small, so things that are distinct in reality get compressed into one point on the mental map. From the inside this does not feel like compression; it feels like seeing one thing. A young child or an ancient Greek philosopher, the author says, would know only one event called ``sound,'' not acoustic vibrations and auditory experiences.

The dangerous case is not two known things under one label but not knowing that there are two things at all, like a detective who must realize that a suspect has an identical twin. Seeing the split is a scientific insight, often marked by new words. ``Consciousness'' covers being awake, reflecting on oneself and subjective experience, which are different puzzles, and confusing them allows bait-and-switch arguments. A thing with contradictory attributes is a good guess for a merged concept, and the rationalist version of Taboo helps to pry it apart.

\responsehead

The post's central distinction is clear and useful: two known things filed under one label are a lesser problem than not knowing that there are two things. The detective who must think of a twin shows it well, and the arithmetic of the map of California is right. The questions concern the post's two main examples, each of which has more history than the post gives it.

The first is sound. ``Sufficiently ancient'' leaves room for earlier thinkers. But the best-known Greek account of sound, Aristotle's \textsc{On the Soul}, made the split the post describes. It explains sound as a movement of air that reaches the ear, and says that ``'sound' and 'hearing' are both ambiguous.'' It also diagnoses as an ambiguity the older claim that there is no colour without sight, a relative of the falling-tree dispute that this sequence uses. This supports the post's thesis. It also shows that the split was made about twenty-three centuries ago.

The second is consciousness. The post's analysis closely matches David Chalmers's 1995 paper. Chalmers calls the word ambiguous, lists ``the difference between wakefulness and sleep'' among the phenomena it covers, takes Crick and Koch's 35 to 75 hertz oscillations as the paper's example of a theory that explains only the easier ones, and calls the resulting move a ``bait-and-switch.'' The post names Chalmers for the hard problem, and a reader who wants the fuller analysis will find it in that paper.

In short: a clear distinction, well illustrated by the detective's twin, whose two main examples have more history than the post gives: Aristotle had already split ``sound,'' and Chalmers's 1995 paper, which the post cites for the hard problem, already analysed ``consciousness'' this way.
\end{afterword}
